\sectionpaged{Zambretti Forecaster}

\subsubsection{Overview}
\marginnote{
    It is reported that the forecasts are most accurate when generated at 9.00 am local time.
}
The original Zambretti Forecaster\index{Zambretti Forecaster}
was developed in 1915 by precision instrument makers Negretti and Zambra.
It was a hand-held disc calculator, that provided on-demand short-term forecasts of up to 12 hours.
It consists of 3 plastic discs of descending radius, arranged concentrically
and secured together by means of a rivet.

The large outside disc A has a short scale calibrated in wind direction,
the middle disc B, has an index to be aligned appropriately with the wind direction scale on A,
and itself carries a scale calibrated in units of barometric pressure (inches of mercury or millibars).

The inner disc C has an index settable to the atmospheric pressure on B and three windows which
allow the selection of a reading based upon of the rising or falling of the barometer in
Winter or Summer.

Through the appropriate window, an alphabetic letter is visible ( printed on disc A)
and this letter is the key to the forecast which can be read off on the obverse side of the instrument.

The calculator uses the mean sea level pressure (MSLP) and pressure trend (steady, rising or falling)
as the primary input.
The season (Summer or Winter) and the wind play a minor role in determining the forecast.

\paragraph{Implementation}

The Zambretti Forecast is implemented with the
\Command{CalculateZambrettiForecast}
function.

This function accepts 5 parameters:

t is the Time of the forecast.

h is the Hemisphere the weather station is located within.

pressure0 is the previous pressure at mean sea level, ideally 3 hours before t.

pressure is the pressure at mean sea level at t.

windDirection is the wind direction at t.

\begin{lstlisting}[caption={Import Zambretti Forecaster},language=Go]
import "github.com/peter-mount/piweather.center/weather/forecast"
\end{lstlisting}
